\section{Odd-$k$ Separation: A General Construction}
\label{sec:odd-k-separation}

The two-point $1$-NN witness of \Cref{sec:finite-separation} shows that LOO
stability and replace-one stability are not interchangeable for $k=1$. A natural
question is whether the same phenomenon persists for larger odd $k$. We answer
this question in the affirmative by constructing, for every odd $k \ge 3$, a
deterministic finite metric sample in which all LOO indicators vanish while a
targeted replace-one perturbation changes the loss at a designated query-label
pair.  (The $k=1$ case is already covered by the two-point witness of
\Cref{sec:finite-separation}.)

\subsection{Construction}

Let $k = 2m+1$ for an integer $m \ge 1$. Let the metric space consist of two
distinct points $a$ and $b$ with $d(a,b) > 0$; any metric space with at least
two points suffices. (For concreteness, one may take the two-vertex graph
metric with $d(a,b)=1$.)

Define the ordered sample
\[
S_m \;=\; \big((a,0)_1,\; \ldots,\; (a,0)_{m+1},\;
          (a,1)_1,\; \ldots,\; (a,1)_{m},\;
          (b,0)_1,\; \ldots,\; (b,0)_{m}\big),
\]
where subscripts track the occurrence index within the ordered tuple.
The sample has size $n = (m+1) + m + m = 3m+1$.

Fix the query-label pair $(x,y) = (a,0)$ and the neighborhood size $k = 2m+1$.

\begin{theorem}[Odd-$k$ LOO/replace-one separation]
\label{thm:odd-k-separation}
For every odd $k = 2m+1$ with $m \ge 1$, the sample $S_m$ satisfies
\[
\Delta_{\mathrm{loo}}(S_m, i) = 0 \qquad \text{for every index } i,
\]
while there exists an index $j$ and a replacement point $z$ such that
\[
\Delta_{\mathrm{rep}}(S_m, j, z, a, 0) = 1.
\]
Thus, pointwise LOO stability does not imply pointwise replace-one stability
for deterministic $k$-NN with odd $k \ge 3$.
\end{theorem}

\begin{proof}
We analyze the margin behavior at the query $a$ and at each deleted occurrence.
Recall from \Cref{prop:5.1} that a deterministic $k$-NN prediction changes
if and only if the signed vote margin $M_k(S, x)$ crosses zero.

\paragraph{Original prediction at $(a,0)$.}
At query $a$, the ordered neighbors are sorted by distance then occurrence index.
All occurrences at $a$ are at distance $0$; all occurrences at $b$ are at
distance $d(a,b) > 0$. The deterministic order therefore selects all
$(m+1)$ copies of $(a,0)$ first, then all $m$ copies of $(a,1)$, giving exactly
$k = 2m+1$ neighbors. The signed margin is
\[
M_k(S_m, a) = (m+1)(-1) + m(+1) = -1 < 0,
\]
so the classifier predicts $0$ at $a$. Hence the loss at $(a,0)$ is $0$.

\paragraph{LOO at an $(a,0)$ occurrence.}
Delete the $i$-th copy of $(a,0)$ and evaluate at the deleted occurrence
$(a,0)$. After deletion the remaining sample has size $n-1 = 3m$, and
$k' = \min(k, n-1) = 2m+1$ for all $m \ge 1$.

At query $a$, the top-$k'$ neighbors consist of the remaining $m$ copies
of $(a,0)$, all $m$ copies of $(a,1)$, and the first copy of $(b,0)$.
The margin is
\[
M_{k'}(S_m^{-i}, a) = m(-1) + m(+1) + 1(-1) = -1.
\]
The prediction is $0$, unchanged from the original. Since the loss at $(a,0)$
was $0$ before and remains $0$ after deletion, $\Delta_{\mathrm{loo}} = 0$.

\paragraph{LOO at an $(a,1)$ occurrence.}
Delete the $i$-th copy of $(a,1)$ and evaluate at $(a,1)$. After deletion,
the top-$k'$ at $a$ consists of all $(m+1)$ copies of $(a,0)$, the remaining
$(m-1)$ copies of $(a,1)$, and the first copy of $(b,0)$. The margin is
\[
M_{k'}(S_m^{-i}, a) = (m+1)(-1) + (m-1)(+1) + 1(-1) = -3.
\]
The prediction is $0$. The original prediction at $(a,1)$ was also $0$
(margin $-1$), so the original loss at $(a,1)$ was $1$; the loss after deletion
is still $1$. Hence $\Delta_{\mathrm{loo}} = 0$.

\paragraph{LOO at a $(b,0)$ occurrence.}
Delete the $i$-th copy of $(b,0)$ and evaluate at $(b,0)$. For the original
sample at query $b$, the $(b,0)$ occurrences are at distance $0$ and appear
first in the top-$k$, followed by the $(a,0)$ occurrences at distance
$d(a,b) > 0$. The original margin is
\[
M_k(S_m, b) = m(-1) + (m+1)(-1) = -(2m+1),
\]
so the original prediction is $0$ and the loss at $(b,0)$ is $0$.

After deletion, the top-$k'$ at $b$ consists of $(m-1)$ copies of $(b,0)$,
all $(m+1)$ copies of $(a,0)$, and the first copy of $(a,1)$. The margin is
\[
M_{k'}(S_m^{-i}, b) = (m-1)(-1) + (m+1)(-1) + 1(+1) = -(2m-1).
\]
For all $m \ge 0$, this is $\le 0$, so the prediction remains $0$ and the
loss stays $0$. Thus $\Delta_{\mathrm{loo}} = 0$.

Since every index falls into one of the three cases above, we have
$\Delta_{\mathrm{loo}}(S_m, i) = 0$ for all $i$.

\paragraph{Replace-one instability.}
Now replace the {\em first} occurrence $(a,0)_1$ with $(a,1)$. The perturbed
sample is
\[
S_m^{1 \leftarrow (a,1)} = \big((a,1),\; (a,0)_2,\ldots,(a,0)_{m+1},\;
                             (a,1)_1,\ldots,(a,1)_m,\;
                             (b,0)_1,\ldots,(b,0)_m\big).
\]

At query $a$, the top-$k$ neighbors are now: $(a,1)$ at index $1$, then
$(a,0)_2$ through $(a,0)_{m+1}$ ($m$ copies), then $(a,1)_1$ through
$(a,1)_m$ ($m$ copies). All $k = 2m+1$ neighbors are at distance $0$. The
signed margin becomes
\[
M_k(S_m^{1 \leftarrow (a,1)}, a) = 1(+1) + m(-1) + m(+1) = 1 > 0,
\]
so the classifier now predicts $1$ at $a$. The loss at $(a,0)$ changes from
$0$ to $1$, giving $\Delta_{\mathrm{rep}}(S_m, 1, (a,1), a, 0) = 1$.
\end{proof}

\paragraph{The case $k=1$.}
The $k=1$ case, excluded from \Cref{thm:odd-k-separation}, is already covered by the two-point witness of \Cref{sec:finite-separation}.  Taken together, the LOO/replace-one separation holds for every odd $k$: $k=1$ by the two-point witness and $k \ge 3$ by \Cref{thm:odd-k-separation}.

\subsection{Remarks on the Construction}

The construction has several noteworthy features.

\paragraph{Role of the metric space.} The argument uses only the fact that
$a$ and $b$ are distinct points with $d(a,b) > 0$; no special distance value,
triangle inequality, or graph structure is required. The separation therefore
holds in {\em any} metric space with at least two points.

\paragraph{Role of sample ordering.} The ordered occurrence sequence is
essential. Placing all $(a,0)$ copies before $(a,1)$ copies ensures that the
original top-$k$ at $a$ contains one more $0$ than $1$, giving margin $-1$.
If the order were reversed, the margin would be $+1$ and the separation would
run in the opposite direction (LOO stable, replace-one from $1$ to $0$).

\paragraph{Role of conflicting labels.} The construction uses conflicting
labels at the same point $a$. This is not an artifact: it is the mechanism
by which a single replacement can change the label composition of the top-$k$
neighborhood without changing its geometric membership. A construction that
forbids conflicting labels would require a more elaborate metric space.

\paragraph{Sample size growth.} The sample size is $3m+1 = \frac{3}{2}(k-1)+1 =
\frac{3k-1}{2}$, which is linear in $k$. The fraction of $(a,0)$ occurrences
among the total is $\frac{m+1}{3m+1} \to \frac{1}{3}$ as $k \to \infty$.

\paragraph{Limit of tie-breaking dependence.} The LOO-stable margin after
deleting an $(a,0)$ occurrence remains $-1$, which is non-zero. This is in
contrast to the even-$k$ case (see \Cref{sec:even-k-analysis}), where the
analogous construction produces margin $0$ after deletion and therefore depends
on tie-breaking to stabilize the prediction.

\subsection{Comparison with Prior Computational Evidence}

The construction above subsumes the computational search evidence for
$k=3,5,7$ that motivated this investigation.  Those candidates were generated
via bounded enumeration and are now superseded by the uniform combinatorial
argument that extends the separation to {\em all} odd $k \ge 3$.
